%% ----------------------------------------------------------------
%% Chapter 4 -- Results and Discussion
%% ----------------------------------------------------------------
\chapter{Results and Discussion}
\label{chap:results}

\section{Experimental Results}
\label{sec:results}

Present your results clearly using figures and tables. Refer to each
figure/table in the text before it appears, e.g.\ \fref{fig:result1}.

\begin{figure}[htbp]
  \centering
  % \includegraphics[width=0.75\textwidth]{Figures/result1.png}
  \caption{Result 1 -- descriptive caption here}
  \label{fig:result1}
\end{figure}

\begin{table}[htbp]
  \centering
  \caption{Performance Comparison}
  \label{tab:comparison}
  \begin{tabular}{llcc}
    \toprule
    \textbf{Method}  & \textbf{Metric 1} & \textbf{Metric 2} & \textbf{Metric 3} \\
    \midrule
    Baseline         & --                & --                & --                \\
    Proposed         & --                & --                & --                \\
    \bottomrule
  \end{tabular}
\end{table}

\subsection{Sub-Result 1}
\label{subsec:result1}

Describe and interpret the first set of results.

\subsection{Sub-Result 2}
\label{subsec:result2}

Describe and interpret the second set of results.

\section{Discussion}
\label{sec:discussion}

Interpret the results in the context of your objectives and the
literature. Address the following:

\begin{itemize}
  \item How do your results compare to existing methods?
  \item Were all objectives met? If not, why?
  \item What are the limitations of your approach?
\end{itemize}

\section{Conclusion}
\label{sec:conclusion}

Summarise the key findings of this project and their significance.
State clearly what was achieved, what was not, and the directions for
future work.
